\documentclass[11pt,a4paper]{article}
\usepackage{curso}
\begin{document}
\portada{Lección}{9.2}{Difusión discreta}

\noindent\textbf{Unidad 9.}\quad\textbf{Objetivos:} difusión enmascarada (estado absorbente); objetivo MDLM; error de revelar tokens en paralelo y muestreadores de pocos pasos.
\textbf{Prerrequisitos:} 4.1, 9.1.

\section{Proceso enmascarado}
Cada token se reemplaza por $[\mathrm M]$ con probabilidad $t\in[0,1]$, independientemente: $q(x_t^i\mid x_0^i)=(1-t)\delta_{x_0^i}+t\,\delta_{[\mathrm M]}$. La máscara es absorbente.
El proceso inverso parte de todo enmascarado y revela tokens con un modelo $p_\theta(x_0^i\mid x_t)$.

\section{Objetivo MDLM}
La cota inferior de la verosimilitud se reduce a una entropía cruzada sobre las posiciones enmascaradas, ponderada en el tiempo (calendario lineal):
\[\mathcal L=\mathbb E_{t,x_t}\Big[\tfrac1t\textstyle\sum_{i:\,x_t^i=[\mathrm M]}-\log p_\theta(x_0^i\mid x_t)\Big].\]
Como BERT, pero sobre todas las tasas de enmascarado y con peso que la convierte en ELBO.

\section{Revelar en paralelo \evS}
Revelar un token por paso es exacto con un modelo perfecto; revelar un grupo $S$ a la vez usa $\prod_{i\in S}p(x^i\mid x_{\text{vis}})\ne p(x_S\mid x_{\text{vis}})$.
Juguete: cadenas binarias de longitud 8 (cadena de Markov que repite el bit anterior con probabilidad 0,9; 48\,\% de cadenas constantes), condicionales exactas,
$2\times10^4$ muestras:
\begin{center}\begin{tabular}{lcccc}\toprule
pasos & 8 & 4 & 2 & 1\\\midrule
variación total & 0,03 & 0,20 & 0,46 & 0,79\\
cadenas constantes & 47\,\% & 29\,\% & 9\,\% & 0,7\,\%\\\bottomrule\end{tabular}\end{center}
El error proviene del muestreo en paralelo, no del modelo. Remedios: revelar por confianza, destilar a pocos pasos, o muestrear la conjunta de cada paso con un EBM
(modelos termodinámicos de eliminación de ruido, 9.3).

\begin{aplicacion}
Modelos de lenguaje por difusión, no autorregresivos; relleno y edición. Laboratorio L13: MDLM pequeño, calidad frente a pasos.
\end{aplicacion}

\begin{error}
\begin{itemize}[nosep]
\item Olvidar que la máscara es absorbente. \quad Quitar el peso $1/t$.
\item Revelar muchos tokens correlacionados y culpar al modelo. \quad Comparar pasos sin igualar evaluaciones.
\end{itemize}
\end{error}

\begin{thebibliography}{9}\small
\bibitem{mdlm} S. S. Sahoo et al., Simple and effective masked diffusion language models, \emph{NeurIPS} (2024).
\bibitem{d3pm} J. Austin et al., Structured denoising diffusion models in discrete state-spaces, \emph{NeurIPS} (2021).
\bibitem[48]{dtm} A. Jelinčič et al., An efficient probabilistic hardware architecture for diffusion-like models, 2025.
\end{thebibliography}
\end{document}
